\begin{frame}{Index 1 of 5: principle, slide, prompt}
\scriptsize
\begin{tabular}{@{}r>{\raggedright\arraybackslash}p{5.4cm}>{\raggedright\arraybackslash}p{4.6cm}p{1.2cm}@{}}
\textbf{Slide} & \textbf{Title} & \textbf{Principle} & \textbf{Prompt}\\\midrule
11 & A relation is a set of pairs & WITNESS / COUNTEREXAMPLE & \\
12 & See it: pairs, matrix, digraph & WITNESS / COUNTEREXAMPLE & 01a\\
14 & Four questions, four witnesses & WITNESS / COUNTEREXAMPLE & 02a\\
15 & Two big kinds of relation & EQUIVALENCE VS PARTIAL ORDER & \\
16 & Sort relations into kinds & EQUIVALENCE VS PARTIAL ORDER & 03a\\
22 & Rules of the game & PARTIAL-ORDER TEST ORDERING & 00a\\
24 & Depends on: the plan as arrows & PARTIAL-ORDER TEST ORDERING & \\
25 & Test order: leaves first & PARTIAL-ORDER TEST ORDERING & \\
26 & Piece 1: letter\_grade & EQUIVALENCE PARTITIONING, BOUNDARIES & \\
\end{tabular}
\note{Use this to find a slide fast: slide number, which lesson idea it shows, and which prompt file it pairs with (prompts/NN, a = first block, b = second block).}
\end{frame}
\begin{frame}{Index 2 of 5: principle, slide, prompt}
\scriptsize
\begin{tabular}{@{}r>{\raggedright\arraybackslash}p{5.4cm}>{\raggedright\arraybackslash}p{4.6cm}p{1.2cm}@{}}
\textbf{Slide} & \textbf{Title} & \textbf{Principle} & \textbf{Prompt}\\\midrule
27 & Piece 1b: ``same bin'' is an equivalence & REFLEXIVE / SYMMETRIC / TRANSITIVE TESTS & \\
28 & Piece 2: similar\_names & EQUIVALENCE PARTITIONING & \\
29 & Piece 3: same\_person must be transitive & TRANSITIVE TEST, WITNESS / COUNTEREXAMPLE & \\
30 & Fix 1: a canonical key & EQUIVALENCE PARTITIONING, WITNESS / COUNTEREXAMPLE & \\
31 & Fix 2: group by chains (the alternative) & EQUIVALENCE PARTITIONING, WITNESS / COUNTEREXAMPLE & \\
32 & Piece 4: is\_duplicate must be symmetric & REFLEXIVE AND SYMMETRIC TESTS & \\
33 & Piece 5: ordering must be antisymmetric & ANTISYMMETRIC TEST, WITNESS / COUNTEREXAMPLE & \\
34 & Piece 6: merge\_duplicates & EQUIVALENCE PARTITIONING & \\
35 & Piece 7: depends-on order & PARTIAL-ORDER TEST ORDERING & \\
\end{tabular}
\note{Use this to find a slide fast: slide number, which lesson idea it shows, and which prompt file it pairs with (prompts/NN, a = first block, b = second block).}
\end{frame}
\begin{frame}{Index 3 of 5: principle, slide, prompt}
\scriptsize
\begin{tabular}{@{}r>{\raggedright\arraybackslash}p{5.4cm}>{\raggedright\arraybackslash}p{4.6cm}p{1.2cm}@{}}
\textbf{Slide} & \textbf{Title} & \textbf{Principle} & \textbf{Prompt}\\\midrule
36 & Piece 8: cycles are a design smell & ANTISYMMETRIC TEST, WITNESS / COUNTEREXAMPLE & \\
38 & Step 1: bins and boundaries, red & EQUIVALENCE PARTITIONING, BOUNDARIES & 04a\\
39 & Step 2: bins and boundaries, green & EQUIVALENCE PARTITIONING, BOUNDARIES & 04b\\
40 & Step 3: the bins are an equivalence & REFLEXIVE / SYMMETRIC / TRANSITIVE TESTS & 05a\\
41 & Step 4: the fuzzy name rule, red & EQUIVALENCE PARTITIONING & 06a\\
42 & Step 5: the fuzzy name rule, green & EQUIVALENCE PARTITIONING & 06b\\
43 & Step 6: the transitivity test, red & TRANSITIVE TEST, WITNESS / COUNTEREXAMPLE & 07a\\
44 & Step 7: fix with a canonical key, green & EQUIVALENCE PARTITIONING, WITNESS / COUNTEREXAMPLE & 08a\\
45 & Step 8: chains (union-find), red & EQUIVALENCE PARTITIONING, WITNESS / COUNTEREXAMPLE & 09a\\
\end{tabular}
\note{Use this to find a slide fast: slide number, which lesson idea it shows, and which prompt file it pairs with (prompts/NN, a = first block, b = second block).}
\end{frame}
\begin{frame}{Index 4 of 5: principle, slide, prompt}
\scriptsize
\begin{tabular}{@{}r>{\raggedright\arraybackslash}p{5.4cm}>{\raggedright\arraybackslash}p{4.6cm}p{1.2cm}@{}}
\textbf{Slide} & \textbf{Title} & \textbf{Principle} & \textbf{Prompt}\\\midrule
46 & Step 9: chains (union-find), green & EQUIVALENCE PARTITIONING, WITNESS / COUNTEREXAMPLE & 09b\\
47 & Step 10: one-sided ``equals'', red & REFLEXIVE AND SYMMETRIC TESTS, WITNESS & 10a\\
48 & Step 11: symmetric equals, green & REFLEXIVE AND SYMMETRIC TESTS, WITNESS & 10b\\
49 & Step 12: sort ties, red & ANTISYMMETRIC TEST, WITNESS / COUNTEREXAMPLE & 11a\\
50 & Step 13: unique sort key, green & ANTISYMMETRIC TEST, WITNESS / COUNTEREXAMPLE & 11b\\
51 & Step 14: merge end to end, red & EQUIVALENCE PARTITIONING & 12a\\
52 & Step 15: merge end to end, green & EQUIVALENCE PARTITIONING & 12b\\
53 & Step 16: depends-on order, red & PARTIAL-ORDER TEST ORDERING & 13a\\
54 & Step 17: depends-on order, green & PARTIAL-ORDER TEST ORDERING & 13b\\
\end{tabular}
\note{Use this to find a slide fast: slide number, which lesson idea it shows, and which prompt file it pairs with (prompts/NN, a = first block, b = second block).}
\end{frame}
\begin{frame}{Index 5 of 5: principle, slide, prompt}
\scriptsize
\begin{tabular}{@{}r>{\raggedright\arraybackslash}p{5.4cm}>{\raggedright\arraybackslash}p{4.6cm}p{1.2cm}@{}}
\textbf{Slide} & \textbf{Title} & \textbf{Principle} & \textbf{Prompt}\\\midrule
55 & Step 18: cycle detection, red & ANTISYMMETRIC TEST, WITNESS / COUNTEREXAMPLE & 14a\\
56 & Step 19: cycle detection, green & ANTISYMMETRIC TEST, WITNESS / COUNTEREXAMPLE & 14b\\
\end{tabular}
\note{Use this to find a slide fast: slide number, which lesson idea it shows, and which prompt file it pairs with (prompts/NN, a = first block, b = second block).}
\end{frame}
